\section{Experimental Configuration Details}
\label{app:experimental-config}

Table~\ref{tab:dataset-spec} reports the generation grid and the partition into task families. The archived field \texttt{n\_if}, displayed as \texttt{if} in the diagnostic exports, corresponds to \texttt{nb\_if} in this paper. All archived configurations record generator mode \texttt{d}, including the zero-condition instances; this historical field does not change the analysis partition. Table~\ref{tab:models} reports the serving configuration recovered from the logs. The logs leave model and tokenizer revisions unset, so the named checkpoints do not establish an immutable weight revision.

\begin{table*}[t]
\centering
\caption{Detailed configuration of the evaluation set. Both task families use the same generation grid.}
\label{tab:dataset-spec}
\scriptsize
\begin{tabularx}{\textwidth}{p{0.22\textwidth}p{0.24\textwidth}X}
\toprule
Category & Parameter & Value \\
\midrule
Dataset size & Java classes & 1,232 per task family \\
& Questions & 46,656 per task family \\
& Question polarity & 23,328 reachable and 23,328 unreachable \\
\midrule
Graph generation & Target class size & $\{100,250,500\}$ methods \\
& Absolute source--target distance & $\{1,4,7,10,13,16,19,22,25\}$ \\
& Local tree structure & Three trees of depth three; two calls per function \\
& Helper-call probability & 0.2 in total \\
\midrule
Generated code & Comments & $\{0,1\}$ \\
& Variables, loops, and parameters & $\{0,1,2\}$ \\
& Conditionals or guards & $\{0,1,3,5\}$ \\
\midrule
Factorial design & Code configurations & $3\times2\times3\times3\times4\times3=648$ \\
& Questions per distance and configuration & Four reachable and four unreachable \\
\midrule
Task semantics & Unguarded & All generated direct calls contribute to reachability \\
& Condition-aware & Only calls whose generated guards evaluate to true contribute to reachability \\
\bottomrule
\end{tabularx}
\end{table*}

\begin{table*}[!t]
\centering
\caption{Model configurations used in the evaluation. Checkpoint and serving information will be added after the final configurations are frozen.}
\label{tab:models}
\scriptsize
\begin{tabularx}{\textwidth}{p{0.25\textwidth}p{0.22\textwidth}p{0.22\textwidth}X}
\toprule
Configuration & Inference mode & Budget metadata & Reproducibility metadata \\
\midrule
\texttt{[MODEL\_CONFIG\_1]} & \texttt{[MODE\_1]} & \texttt{[BUDGET\_CONTROL\_1]} & checkpoint, size, context, license and source, serving settings \\
\texttt{[MODEL\_CONFIG\_2]} & \texttt{[MODE\_2]} & \texttt{[BUDGET\_CONTROL\_2]} & checkpoint, size, context, license and source, serving settings \\
\texttt{[MODEL\_CONFIG\_3]} & \texttt{[MODE\_3]} & \texttt{[BUDGET\_CONTROL\_3]} & checkpoint, size, context, license and source, serving settings \\
\texttt{[MODEL\_CONFIG\_4]} & \texttt{[MODE\_4]} & \texttt{[BUDGET\_CONTROL\_4]} & checkpoint, size, context, license and source, serving settings \\
\texttt{[MODEL\_CONFIG\_5]} & \texttt{[MODE\_5]} & \texttt{[BUDGET\_CONTROL\_5]} & checkpoint, size, context, license and source, serving settings \\
\bottomrule
\end{tabularx}
\end{table*}


\section{Diagnostic Definitions and Populations}
\label{app:diagnostic-definitions}

Let $R(s)$ be the set of methods executable from $s$, and $T$ the returned sequence. A guard is exposed when its containing method appears in $T$. A false-guard error is detected when a newly introduced method has no executable predecessor among earlier trace methods and exactly one earlier guard explains its introduction. Multiple candidate guards make attribution undefined. This diagnostic identifies unsupported inclusions through false guards; it does not observe all possible errors of condition evaluation, such as ignoring a true branch.

\vspace{0.2cm}

T01 is the number of attributed guard errors divided by the number of exposed call-controlling guards, summed over eligible responses. A is the fraction of attributed errors whose guarded callees can reach the target through executable edges. B is the fraction of eligible responses containing at least one such error. T04 is the mean fraction of $R(s)$ omitted on negative questions; T05 is the incidence of a negative trace containing an existing method outside $R(s)$. They measure different, potentially overlapping defects.

\vspace{0.2cm}

T06 measures target omission among structurally admissible positive traces; T07 measures continuation after the target among such traces containing it. Structural admissibility requires a source-first sequence with no duplicate or unsupported method introduction. T08 measures unknown method names. T09 reports the relative position of the first structural error among nonempty defective traces. T11 measures duplicate methods in eligible nonempty traces. T13 is mean source-anchored coverage: the fraction of $R(s)$ linked back to the source through executable introductions. It is not a positive-query success score, since a positive trace can correctly stop at the target before covering the whole region.

\vspace{0.2cm}

For a guard $g$ with resolved variable origin, let $h(g)$ count method boundaries back to the local declaration supplying its tested value. D1 is the mean $h(g)$ over exposed guards; D3 restricts to guards with an attributed false-guard error, and D5 to the remaining exposed guards in unambiguous traces. Literal conditions, unresolved origins, and ambiguous lineages are excluded. D1$^*$ additionally uses the attribution population shared by D3/D5, so its numerator and denominator partition into those of D3 and D5. D9 first averages $h(g)$ over guards in each reference depth-first exploration, then averages the resulting values across questions. The reference traversal is a diagnostic convention, not a mandatory model traversal.

\vspace{0.2cm}

Table~\ref{tab:difficulty-populations} makes these different weights and exclusions explicit. The unrestricted D1 differs from D1$^*$ because it also includes resolved exposures in traces with ambiguous error attribution. Table~\ref{tab:diagnostic-populations} retains the populations of the other reported diagnostics. Missing measurements are not interpreted as successful evaluations. The source export and its metric catalogue are identified by the hashes accompanying the reproducible aggregates.

% Generated by python3 scripts/render_results.py; source: data/results/aggregates.csv.
\begin{table*}[t]
\centering
\caption{Difficulty populations within condition-aware questions. Each parameter level has 11,664 questions. The denominator counts resolved guard occurrences for D1/D1$^*$/D3/D5, and eligible questions for D9. Exclusions include zero resolved occurrences, missing origins, or ambiguous attribution as applicable.}
\label{tab:difficulty-populations}
\scriptsize
\setlength{\tabcolsep}{3pt}
\begin{tabular}{llrrrrr}
\toprule
Model & Param. & Metric & Mean & Eligible questions & Denominator & Excluded \\
\midrule
Gemma & 1 & D1 & 3.984 & 11,149 & 161,390 & 515 \\
Gemma & 1 & D1$^*$ & 3.724 & 9,337 & 95,113 & 2,327 \\
Gemma & 1 & D3 & 5.411 & 2,784 & 7,057 & 8,880 \\
Gemma & 1 & D5 & 3.588 & 9,328 & 88,056 & 2,336 \\
Gemma & 1 & D9 & 4.251 & 11,592 & 11,592 & 72 \\
Gemma & 2 & D1 & 4.922 & 11,073 & 163,650 & 591 \\
Gemma & 2 & D1$^*$ & 4.653 & 8,837 & 83,662 & 2,827 \\
Gemma & 2 & D3 & 6.223 & 3,202 & 8,701 & 8,462 \\
Gemma & 2 & D5 & 4.471 & 8,821 & 74,961 & 2,843 \\
Gemma & 2 & D9 & 5.285 & 11,575 & 11,575 & 89 \\
\midrule
Qwen & 1 & D1 & 3.928 & 11,353 & 143,319 & 311 \\
Qwen & 1 & D1$^*$ & 3.720 & 9,716 & 91,240 & 1,948 \\
Qwen & 1 & D3 & 5.145 & 3,279 & 8,222 & 8,385 \\
Qwen & 1 & D5 & 3.579 & 9,707 & 83,018 & 1,957 \\
Qwen & 1 & D9 & 4.251 & 11,592 & 11,592 & 72 \\
Qwen & 2 & D1 & 4.789 & 11,342 & 145,326 & 322 \\
Qwen & 2 & D1$^*$ & 4.625 & 9,400 & 84,280 & 2,264 \\
Qwen & 2 & D3 & 5.898 & 3,604 & 9,366 & 8,060 \\
Qwen & 2 & D5 & 4.466 & 9,382 & 74,914 & 2,282 \\
Qwen & 2 & D9 & 5.285 & 11,575 & 11,575 & 89 \\
\bottomrule
\end{tabular}

\end{table*}

% Generated by python3 scripts/render_results.py; source: data/results/aggregates.csv.
\begin{table*}[t]
\centering
\caption{Condition-aware diagnostic aggregates. Values are percentages; T09 is a median with no sum denominator. Eligibility depends on the metric, so exclusions also include the other polarity for negative/positive-only diagnostics. For T04/T13, the numerator sums per-question fractions, not node counts. The two FP rows use unambiguously attributable negative records and condition on whether a target-directed guard error is present.}
\label{tab:diagnostic-populations}
\scriptsize
\setlength{\tabcolsep}{3pt}
\begin{tabular}{llrrrrr}
\toprule
Model & Metric & Value & Numerator & Denominator & Eligible & Excluded \\
\midrule
Gemma & T01 & 4.34 & 16,253.00 & 374,856 & 29,507 & 5,485 \\
Gemma & T04 & 39.37 & 6,888.62 & 17,496 & 17,496 & 17,496 \\
Gemma & T05 & 51.38 & 8,990.00 & 17,496 & 17,496 & 17,496 \\
Gemma & T06 & 8.48 & 1,135.00 & 13,392 & 13,392 & 21,600 \\
Gemma & T07 & 1.13 & 138.00 & 12,257 & 12,257 & 22,735 \\
Gemma & T08 & 0.24 & 85.00 & 34,992 & 34,992 & 0 \\
Gemma & T09 & 28.57 & -- & -- & 13,043 & 21,949 \\
Gemma & T11 & 5.71 & 1,943.00 & 34,014 & 34,014 & 978 \\
Gemma & T13 & 44.98 & 15,289.65 & 33,990 & 33,990 & 1,002 \\
Gemma & A & 19.81 & 3,220.00 & 16,253 & 30,784 & 4,208 \\
Gemma & B & 9.65 & 2,971.00 & 30,784 & 30,784 & 4,208 \\
Gemma & FP with target-directed error & 89.73 & 1,721.00 & 1,918 & 1,918 & 33,074 \\
Gemma & FP without target-directed error & 1.36 & 170.00 & 12,536 & 12,536 & 22,456 \\
\midrule
Qwen & T01 & 5.34 & 17,874.00 & 334,698 & 30,526 & 4,466 \\
Qwen & T04 & 38.01 & 6,650.81 & 17,496 & 17,496 & 17,496 \\
Qwen & T05 & 48.50 & 8,485.00 & 17,496 & 17,496 & 17,496 \\
Qwen & T06 & 4.12 & 614.00 & 14,893 & 14,893 & 20,099 \\
Qwen & T07 & 0.06 & 8.00 & 14,279 & 14,279 & 20,713 \\
Qwen & T08 & 0.59 & 207.00 & 34,992 & 34,992 & 0 \\
Qwen & T09 & 30.00 & -- & -- & 11,227 & 23,765 \\
Qwen & T11 & 1.29 & 445.00 & 34,534 & 34,534 & 458 \\
Qwen & T13 & 41.59 & 14,296.55 & 34,377 & 34,377 & 615 \\
Qwen & A & 23.30 & 4,165.00 & 17,874 & 31,291 & 3,701 \\
Qwen & B & 12.22 & 3,824.00 & 31,291 & 31,291 & 3,701 \\
Qwen & FP with target-directed error & 92.57 & 2,367.00 & 2,557 & 2,557 & 32,435 \\
Qwen & FP without target-directed error & 0.46 & 55.00 & 12,074 & 12,074 & 22,918 \\
\bottomrule
\end{tabular}

\end{table*}


\section{Detailed Result Curves}
\label{app:result-details}

Figure~\ref{fig:results-guards-parameters} shows the marginal guard and parameter profiles. Table~\ref{tab:guard-mechanisms} distinguishes the rate of detected guard errors from their target-directed consequences. T01 decreases as the guard count rises, even while path scores deteriorate and A increases. More guards increase the exposure denominator, and the returned traces change which guards are exposed. A smaller T01 ratio is therefore not evidence that the model has become better at condition evaluation. Counts and eligibility must be considered together.

\vspace{0.2cm}

Figure~\ref{fig:results-distance} separates signed distance levels rather than pooling polarities. Figure~\ref{fig:results-context} shows that higher negative answer correctness can coexist with lower path correctness. These curves complement the primary scores by locating where the decision and its evidence diverge.

% Generated by python3 scripts/render_results.py; source: data/results/aggregates.csv.
\begin{figure*}[t]
\centering
\begin{tikzpicture}
\begin{groupplot}[rb score axis, group style={group size=2 by 1, horizontal sep=1.3cm}]
\nextgroupplot[title={(a) Guard count},xmin=-0.2,xmax=5.2,xtick={0,1,3,5},xlabel={Generated guards ($\mathtt{nb\_if}$)},legend to name=fig-results-guards-parameters-legend,legend columns=4]
\addplot[rb gemma,rb answer] coordinates {(0,96.0477) (1,89.5319) (3,78.3779) (5,72.8738)};
\addlegendentry{Gemma: answer}
\addplot[rb gemma,rb path] coordinates {(0,82.2274) (1,59.3879) (3,44.5473) (5,37.4486)};
\addlegendentry{Gemma: trace}
\addplot[rb qwen,rb answer] coordinates {(0,99.1941) (1,92.5840) (3,82.3045) (5,77.8464)};
\addlegendentry{Qwen: answer}
\addplot[rb qwen,rb path] coordinates {(0,96.3306) (1,67.7726) (3,57.0302) (5,48.6025)};
\addlegendentry{Qwen: trace}
\nextgroupplot[title={(b) Parameter count},xmin=-0.1,xmax=2.1,xtick={0,1,2},xlabel={Parameters per method},ylabel={}]
\addplot[rb gemma,rb answer] coordinates {(0,91.2380) (1,75.7888) (2,73.7569)};
\addplot[rb gemma,rb path] coordinates {(0,66.6152) (1,39.1718) (2,35.5967)};
\addplot[rb qwen,rb answer] coordinates {(0,97.4108) (1,78.4379) (2,76.8861)};
\addplot[rb qwen,rb path] coordinates {(0,90.5350) (1,43.5185) (2,39.3519)};
\end{groupplot}
\node[anchor=north,inner sep=0pt] at ($(group c1r1.south)!0.5!(group c2r1.south)+(0,-1.05cm)$) {\pgfplotslegendfromname{fig-results-guards-parameters-legend}};
\end{tikzpicture}
\caption{Scores by guard count over all questions (a), and by parameter count within condition-aware questions (b). Other factors are balanced. The zero-guard level is the unguarded group.}
\Description{Two panels with correctness on the same zero-to-one-hundred percent scale. Circles denote Gemma and squares Qwen; solid lines denote answer correctness and dashed lines trace correctness.}
\label{fig:results-guards-parameters}
\end{figure*}

% Generated by python3 scripts/render_results.py; source: data/results/aggregates.csv.
\begin{table*}[t]
\centering
\caption{Guard diagnostics (\%). A is the target-directed share of attributed errors; B is the incidence of at least one such error in an eligible response. T05 applies only to the 5,832 negative questions per guard level. T01 uses the exposed-guard denominator; A/B exclude ambiguous attributions.}
\label{tab:guard-mechanisms}
\scriptsize
\setlength{\tabcolsep}{3pt}
\begin{tabular}{llrrrrrr}
\toprule
Model & Guards & Answer & Path & T01 & A & B & T05 \\
\midrule
Gemma & 1 & 89.53 & 59.39 & 9.21 & 9.24 & 6.07 & 49.35 \\
Gemma & 3 & 78.38 & 44.55 & 3.81 & 25.91 & 11.92 & 52.52 \\
Gemma & 5 & 72.87 & 37.45 & 2.07 & 36.41 & 11.62 & 52.28 \\
Qwen & 1 & 92.58 & 67.77 & 10.88 & 11.19 & 7.34 & 42.37 \\
Qwen & 3 & 82.30 & 57.03 & 4.70 & 31.01 & 15.10 & 49.49 \\
Qwen & 5 & 77.85 & 48.60 & 2.94 & 36.63 & 15.06 & 53.64 \\
\bottomrule
\end{tabular}

\end{table*}

% Generated by python3 scripts/render_results.py; source: data/results/aggregates.csv.
\begin{figure*}[t]
\centering
\begin{tikzpicture}[x=1cm,y=0.035cm]
\node[font=\small] at (3.4,113) {Final-answer correctness};
\draw[gray!20] (0,0) -- (6.8,0);
\node[anchor=east,font=\scriptsize] at (-0.08,0) {0};
\draw[gray!20] (0,25) -- (6.8,25);
\node[anchor=east,font=\scriptsize] at (-0.08,25) {25};
\draw[gray!20] (0,50) -- (6.8,50);
\node[anchor=east,font=\scriptsize] at (-0.08,50) {50};
\draw[gray!20] (0,75) -- (6.8,75);
\node[anchor=east,font=\scriptsize] at (-0.08,75) {75};
\draw[gray!20] (0,100) -- (6.8,100);
\node[anchor=east,font=\scriptsize] at (-0.08,100) {100};
\draw[gray!60] (0,100) -- (0,0) -- (6.8,0);
\node[anchor=north,font=\scriptsize] at (0.000,-2) {-25};
\node[anchor=north,font=\scriptsize] at (0.408,-2) {};
\node[anchor=north,font=\scriptsize] at (0.816,-2) {-19};
\node[anchor=north,font=\scriptsize] at (1.224,-2) {};
\node[anchor=north,font=\scriptsize] at (1.632,-2) {-13};
\node[anchor=north,font=\scriptsize] at (2.040,-2) {};
\node[anchor=north,font=\scriptsize] at (2.448,-2) {-7};
\node[anchor=north,font=\scriptsize] at (2.856,-2) {};
\node[anchor=north,font=\scriptsize] at (3.264,-2) {-1};
\node[anchor=north,font=\scriptsize] at (3.536,-2) {};
\node[anchor=north,font=\scriptsize] at (3.944,-2) {4};
\node[anchor=north,font=\scriptsize] at (4.352,-2) {};
\node[anchor=north,font=\scriptsize] at (4.760,-2) {10};
\node[anchor=north,font=\scriptsize] at (5.168,-2) {};
\node[anchor=north,font=\scriptsize] at (5.576,-2) {16};
\node[anchor=north,font=\scriptsize] at (5.984,-2) {};
\node[anchor=north,font=\scriptsize] at (6.392,-2) {22};
\node[anchor=north,font=\scriptsize] at (6.800,-2) {25};
\draw[gray!50,dotted] (3.400,0) -- (3.400,100);
\draw[modelorange,line width=0.9pt] plot coordinates {(0.000,76.0802) (0.408,76.5947) (0.816,77.1605) (1.224,77.4691) (1.632,76.8519) (2.040,77.8292) (2.448,78.1379) (2.856,79.0123) (3.264,83.7449) (3.536,97.5309) (3.944,89.0432) (4.352,85.3909) (4.760,82.6132) (5.168,80.8642) (5.576,78.6008) (5.984,77.2119) (6.392,74.8971) (6.800,75.6687)};
\fill[modelorange] (0.000,76.0802) circle (1.2pt);
\fill[modelorange] (0.408,76.5947) circle (1.2pt);
\fill[modelorange] (0.816,77.1605) circle (1.2pt);
\fill[modelorange] (1.224,77.4691) circle (1.2pt);
\fill[modelorange] (1.632,76.8519) circle (1.2pt);
\fill[modelorange] (2.040,77.8292) circle (1.2pt);
\fill[modelorange] (2.448,78.1379) circle (1.2pt);
\fill[modelorange] (2.856,79.0123) circle (1.2pt);
\fill[modelorange] (3.264,83.7449) circle (1.2pt);
\fill[modelorange] (3.536,97.5309) circle (1.2pt);
\fill[modelorange] (3.944,89.0432) circle (1.2pt);
\fill[modelorange] (4.352,85.3909) circle (1.2pt);
\fill[modelorange] (4.760,82.6132) circle (1.2pt);
\fill[modelorange] (5.168,80.8642) circle (1.2pt);
\fill[modelorange] (5.576,78.6008) circle (1.2pt);
\fill[modelorange] (5.984,77.2119) circle (1.2pt);
\fill[modelorange] (6.392,74.8971) circle (1.2pt);
\fill[modelorange] (6.800,75.6687) circle (1.2pt);
\draw[modelblue,line width=0.9pt] plot coordinates {(0.000,73.6111) (0.408,75.6687) (0.816,74.6399) (1.224,74.2798) (1.632,75.0000) (2.040,74.1255) (2.448,77.3148) (2.856,76.1317) (3.264,81.8930) (3.536,99.9486) (3.944,97.8909) (4.352,95.7305) (4.760,93.5185) (5.168,92.5412) (5.576,89.6091) (5.984,88.6317) (6.392,88.3745) (6.800,87.5000)};
\fill[modelblue] (0.000,73.6111) circle (1.2pt);
\fill[modelblue] (0.408,75.6687) circle (1.2pt);
\fill[modelblue] (0.816,74.6399) circle (1.2pt);
\fill[modelblue] (1.224,74.2798) circle (1.2pt);
\fill[modelblue] (1.632,75.0000) circle (1.2pt);
\fill[modelblue] (2.040,74.1255) circle (1.2pt);
\fill[modelblue] (2.448,77.3148) circle (1.2pt);
\fill[modelblue] (2.856,76.1317) circle (1.2pt);
\fill[modelblue] (3.264,81.8930) circle (1.2pt);
\fill[modelblue] (3.536,99.9486) circle (1.2pt);
\fill[modelblue] (3.944,97.8909) circle (1.2pt);
\fill[modelblue] (4.352,95.7305) circle (1.2pt);
\fill[modelblue] (4.760,93.5185) circle (1.2pt);
\fill[modelblue] (5.168,92.5412) circle (1.2pt);
\fill[modelblue] (5.576,89.6091) circle (1.2pt);
\fill[modelblue] (5.984,88.6317) circle (1.2pt);
\fill[modelblue] (6.392,88.3745) circle (1.2pt);
\fill[modelblue] (6.800,87.5000) circle (1.2pt);
\end{tikzpicture}
\hfill
\begin{tikzpicture}[x=1cm,y=0.035cm]
\node[font=\small] at (3.4,113) {Path correctness};
\draw[gray!20] (0,0) -- (6.8,0);
\node[anchor=east,font=\scriptsize] at (-0.08,0) {0};
\draw[gray!20] (0,25) -- (6.8,25);
\node[anchor=east,font=\scriptsize] at (-0.08,25) {25};
\draw[gray!20] (0,50) -- (6.8,50);
\node[anchor=east,font=\scriptsize] at (-0.08,50) {50};
\draw[gray!20] (0,75) -- (6.8,75);
\node[anchor=east,font=\scriptsize] at (-0.08,75) {75};
\draw[gray!20] (0,100) -- (6.8,100);
\node[anchor=east,font=\scriptsize] at (-0.08,100) {100};
\draw[gray!60] (0,100) -- (0,0) -- (6.8,0);
\node[anchor=north,font=\scriptsize] at (0.000,-2) {-25};
\node[anchor=north,font=\scriptsize] at (0.408,-2) {};
\node[anchor=north,font=\scriptsize] at (0.816,-2) {-19};
\node[anchor=north,font=\scriptsize] at (1.224,-2) {};
\node[anchor=north,font=\scriptsize] at (1.632,-2) {-13};
\node[anchor=north,font=\scriptsize] at (2.040,-2) {};
\node[anchor=north,font=\scriptsize] at (2.448,-2) {-7};
\node[anchor=north,font=\scriptsize] at (2.856,-2) {};
\node[anchor=north,font=\scriptsize] at (3.264,-2) {-1};
\node[anchor=north,font=\scriptsize] at (3.536,-2) {};
\node[anchor=north,font=\scriptsize] at (3.944,-2) {4};
\node[anchor=north,font=\scriptsize] at (4.352,-2) {};
\node[anchor=north,font=\scriptsize] at (4.760,-2) {10};
\node[anchor=north,font=\scriptsize] at (5.168,-2) {};
\node[anchor=north,font=\scriptsize] at (5.576,-2) {16};
\node[anchor=north,font=\scriptsize] at (5.984,-2) {};
\node[anchor=north,font=\scriptsize] at (6.392,-2) {22};
\node[anchor=north,font=\scriptsize] at (6.800,-2) {25};
\draw[gray!50,dotted] (3.400,0) -- (3.400,100);
\draw[modelorange,line width=0.9pt,dashed] plot coordinates {(0.000,11.6255) (0.408,14.1461) (0.816,16.0494) (1.224,19.5988) (1.632,20.8333) (2.040,23.9198) (2.448,28.4465) (2.856,36.2654) (3.264,54.0123) (3.536,95.6790) (3.944,81.9959) (4.352,72.9938) (4.760,69.2387) (5.168,65.5350) (5.576,64.0432) (5.984,59.9794) (6.392,57.3560) (6.800,56.5844)};
\fill[modelorange] (0.000,11.6255) circle (1.2pt);
\fill[modelorange] (0.408,14.1461) circle (1.2pt);
\fill[modelorange] (0.816,16.0494) circle (1.2pt);
\fill[modelorange] (1.224,19.5988) circle (1.2pt);
\fill[modelorange] (1.632,20.8333) circle (1.2pt);
\fill[modelorange] (2.040,23.9198) circle (1.2pt);
\fill[modelorange] (2.448,28.4465) circle (1.2pt);
\fill[modelorange] (2.856,36.2654) circle (1.2pt);
\fill[modelorange] (3.264,54.0123) circle (1.2pt);
\fill[modelorange] (3.536,95.6790) circle (1.2pt);
\fill[modelorange] (3.944,81.9959) circle (1.2pt);
\fill[modelorange] (4.352,72.9938) circle (1.2pt);
\fill[modelorange] (4.760,69.2387) circle (1.2pt);
\fill[modelorange] (5.168,65.5350) circle (1.2pt);
\fill[modelorange] (5.576,64.0432) circle (1.2pt);
\fill[modelorange] (5.984,59.9794) circle (1.2pt);
\fill[modelorange] (6.392,57.3560) circle (1.2pt);
\fill[modelorange] (6.800,56.5844) circle (1.2pt);
\draw[modelblue,line width=0.9pt,dashed] plot coordinates {(0.000,22.9424) (0.408,25.6173) (0.816,27.5720) (1.224,30.0926) (1.632,31.5844) (2.040,34.2593) (2.448,37.9115) (2.856,40.7922) (3.264,55.5556) (3.536,99.8971) (3.944,91.4095) (4.352,86.6770) (4.760,81.0185) (5.168,81.0700) (5.576,75.9259) (5.984,75.6687) (6.392,71.9136) (6.800,70.5247)};
\fill[modelblue] (0.000,22.9424) circle (1.2pt);
\fill[modelblue] (0.408,25.6173) circle (1.2pt);
\fill[modelblue] (0.816,27.5720) circle (1.2pt);
\fill[modelblue] (1.224,30.0926) circle (1.2pt);
\fill[modelblue] (1.632,31.5844) circle (1.2pt);
\fill[modelblue] (2.040,34.2593) circle (1.2pt);
\fill[modelblue] (2.448,37.9115) circle (1.2pt);
\fill[modelblue] (2.856,40.7922) circle (1.2pt);
\fill[modelblue] (3.264,55.5556) circle (1.2pt);
\fill[modelblue] (3.536,99.8971) circle (1.2pt);
\fill[modelblue] (3.944,91.4095) circle (1.2pt);
\fill[modelblue] (4.352,86.6770) circle (1.2pt);
\fill[modelblue] (4.760,81.0185) circle (1.2pt);
\fill[modelblue] (5.168,81.0700) circle (1.2pt);
\fill[modelblue] (5.576,75.9259) circle (1.2pt);
\fill[modelblue] (5.984,75.6687) circle (1.2pt);
\fill[modelblue] (6.392,71.9136) circle (1.2pt);
\fill[modelblue] (6.800,70.5247) circle (1.2pt);
\end{tikzpicture}
\par\smallskip\scriptsize Qwen: blue; Gemma: orange. Solid: answer; dashed: path.
\caption{Condition-aware scores (\%) by signed generation distance. Negative distances denote unreachable targets. Each point has 1,944 questions per model. The dotted line separates polarities; there is no zero-distance query.}
\label{fig:results-distance}
\end{figure*}

% Generated by python3 scripts/render_results.py; source: data/results/aggregates.csv.
\begin{figure*}[t]
\centering
\begin{tikzpicture}[x=1cm,y=0.035cm]
\node[font=\small] at (3.4,113) {Reachable targets};
\draw[gray!20] (0,0) -- (6.8,0);
\node[anchor=east,font=\scriptsize] at (-0.08,0) {0};
\draw[gray!20] (0,25) -- (6.8,25);
\node[anchor=east,font=\scriptsize] at (-0.08,25) {25};
\draw[gray!20] (0,50) -- (6.8,50);
\node[anchor=east,font=\scriptsize] at (-0.08,50) {50};
\draw[gray!20] (0,75) -- (6.8,75);
\node[anchor=east,font=\scriptsize] at (-0.08,75) {75};
\draw[gray!20] (0,100) -- (6.8,100);
\node[anchor=east,font=\scriptsize] at (-0.08,100) {100};
\draw[gray!60] (0,100) -- (0,0) -- (6.8,0);
\node[anchor=north,font=\scriptsize] at (0.000,-2) {100};
\node[anchor=north,font=\scriptsize] at (2.550,-2) {250};
\node[anchor=north,font=\scriptsize] at (6.800,-2) {500};
\draw[modelorange,line width=0.9pt] plot coordinates {(0.000,96.0391) (2.550,85.2023) (6.800,66.0322)};
\fill[modelorange] (0.000,96.0391) circle (1.2pt);
\fill[modelorange] (2.550,85.2023) circle (1.2pt);
\fill[modelorange] (6.800,66.0322) circle (1.2pt);
\draw[modelorange,line width=0.9pt,dashed] plot coordinates {(0.000,79.1324) (2.550,71.6049) (6.800,57.0645)};
\fill[modelorange] (0.000,79.1324) circle (1.2pt);
\fill[modelorange] (2.550,71.6049) circle (1.2pt);
\fill[modelorange] (6.800,57.0645) circle (1.2pt);
\draw[modelblue,line width=0.9pt] plot coordinates {(0.000,97.6680) (2.550,93.5700) (6.800,86.6770)};
\fill[modelblue] (0.000,97.6680) circle (1.2pt);
\fill[modelblue] (2.550,93.5700) circle (1.2pt);
\fill[modelblue] (6.800,86.6770) circle (1.2pt);
\draw[modelblue,line width=0.9pt,dashed] plot coordinates {(0.000,83.8649) (2.550,82.5789) (6.800,78.2579)};
\fill[modelblue] (0.000,83.8649) circle (1.2pt);
\fill[modelblue] (2.550,82.5789) circle (1.2pt);
\fill[modelblue] (6.800,78.2579) circle (1.2pt);
\end{tikzpicture}
\hfill
\begin{tikzpicture}[x=1cm,y=0.035cm]
\node[font=\small] at (3.4,113) {Unreachable targets};
\draw[gray!20] (0,0) -- (6.8,0);
\node[anchor=east,font=\scriptsize] at (-0.08,0) {0};
\draw[gray!20] (0,25) -- (6.8,25);
\node[anchor=east,font=\scriptsize] at (-0.08,25) {25};
\draw[gray!20] (0,50) -- (6.8,50);
\node[anchor=east,font=\scriptsize] at (-0.08,50) {50};
\draw[gray!20] (0,75) -- (6.8,75);
\node[anchor=east,font=\scriptsize] at (-0.08,75) {75};
\draw[gray!20] (0,100) -- (6.8,100);
\node[anchor=east,font=\scriptsize] at (-0.08,100) {100};
\draw[gray!60] (0,100) -- (0,0) -- (6.8,0);
\node[anchor=north,font=\scriptsize] at (0.000,-2) {100};
\node[anchor=north,font=\scriptsize] at (2.550,-2) {250};
\node[anchor=north,font=\scriptsize] at (6.800,-2) {500};
\draw[modelorange,line width=0.9pt] plot coordinates {(0.000,69.2387) (2.550,77.4177) (6.800,87.6372)};
\fill[modelorange] (0.000,69.2387) circle (1.2pt);
\fill[modelorange] (2.550,77.4177) circle (1.2pt);
\fill[modelorange] (6.800,87.6372) circle (1.2pt);
\draw[modelorange,line width=0.9pt,dashed] plot coordinates {(0.000,32.2702) (2.550,25.3086) (6.800,17.3868)};
\fill[modelorange] (0.000,32.2702) circle (1.2pt);
\fill[modelorange] (2.550,25.3086) circle (1.2pt);
\fill[modelorange] (6.800,17.3868) circle (1.2pt);
\draw[modelblue,line width=0.9pt] plot coordinates {(0.000,67.2497) (2.550,75.6344) (6.800,84.6708)};
\fill[modelblue] (0.000,67.2497) circle (1.2pt);
\fill[modelblue] (2.550,75.6344) circle (1.2pt);
\fill[modelblue] (6.800,84.6708) circle (1.2pt);
\draw[modelblue,line width=0.9pt,dashed] plot coordinates {(0.000,36.3512) (2.550,34.0364) (6.800,31.7215)};
\fill[modelblue] (0.000,36.3512) circle (1.2pt);
\fill[modelblue] (2.550,34.0364) circle (1.2pt);
\fill[modelblue] (6.800,31.7215) circle (1.2pt);
\end{tikzpicture}
\par\smallskip\scriptsize Qwen: blue; Gemma: orange. Solid: answer; dashed: path.
\caption{Condition-aware scores (\%) by class size, separated by polarity. Each point has 5,832 questions per model. Negative answer correctness rises while negative path correctness falls.}
\label{fig:results-context}
\end{figure*}


\vspace{0.2cm}

Table~\ref{tab:token-observability} reports persisted token counts. Condition-aware responses contain more reasoning tokens on average than unguarded responses, while their primary scores are lower. Mean reasoning usage is 45.72\% of the requested budget for Gemma and 38.43\% for Qwen on condition-aware questions. These means do not determine whether individual failures exhausted their budget, and tokenizer differences prevent interpreting a token ratio as equal compute. With one requested budget, the experiment does not establish whether increasing it would resolve the observed failures.

% Generated by python3 scripts/render_results.py; source: data/results/aggregates.csv.
\begin{table*}[t]
\centering
\caption{Mean persisted reasoning and answer tokens, counted with each model\'s tokenizer. Usage is reasoning tokens divided by the requested 16,384-token budget. These are output measurements, not wall-clock compute or a budget intervention.}
\label{tab:token-observability}
\scriptsize
\setlength{\tabcolsep}{3pt}
\begin{tabular}{llrrr}
\toprule
Model & Group & Reasoning tokens & Usage (\%) & Answer tokens \\
\midrule
Gemma & Unguarded & 4,345 & 26.52 & 396 \\
Gemma & Condition-aware & 7,490 & 45.72 & 585 \\
Qwen & Unguarded & 3,388 & 20.68 & 83 \\
Qwen & Condition-aware & 6,296 & 38.43 & 96 \\
\bottomrule
\end{tabular}

\end{table*}


\section{Published Benchmark Comparisons}
\label{app:external-comparison}

Table~\ref{tab:external-comparison} records scores from official model cards for the named checkpoints, checked on October~6, 2026. Qwen's card is the source of both model columns except that Google's card independently confirms Gemma's LiveCodeBench score~\cite{qwen2026qwen36,google2026gemma4}. These are reported external results, not additional runs in our experiment.

\begin{table*}[t]
\centering
\caption{Published external scores (\%) for the named checkpoints. Comparisons are within rows; protocols differ across rows and from \benchmark.}
\label{tab:external-comparison}
\scriptsize
\begin{tabularx}{\textwidth}{lrrX}
\toprule
Benchmark & Gemma & Qwen & Provenance and protocol qualification \\
\midrule
LiveCodeBench v6 & 77.1 & 80.4 & Qwen card; Gemma value also confirmed by Google. \\
SWE-bench Verified & 17.4 & 73.4 & Qwen card; internal agent scaffold with shell and file-edit tools. \\
SWE-bench Multilingual & 17.3 & 67.2 & Same vendor-reported scaffold, not our reachability protocol. \\
SWE-bench Pro & 13.8 & 49.5 & Qwen card; refined benchmark set applied to its baselines. \\
Terminal-Bench 2.0 & 34.2 & 51.5 & Qwen card; Harbor/Terminus-2 harness, mean of five runs. \\
\bottomrule
\end{tabularx}
\end{table*}

% Flush appendix floats and restore full-height columns before references.
\clearpage
\nolinenumbers
\twocolumn
